\documentclass[12pt,a4paper]{article}
% \usepackage[utf8]{inputenc} % fontspecと併用時は不要
\usepackage[japanese]{babel}
\usepackage{amsmath}
\usepackage{amsfonts}
\usepackage{amssymb}
\usepackage{geometry}
\usepackage{graphicx}
\usepackage{enumitem}
\usepackage{fancyhdr}
\usepackage{verbatim}

% 日本語改行設定
\usepackage{xeCJK}
\usepackage{indentfirst}
\setlength{\parindent}{1em}

% 日本語フォント設定
\usepackage{fontspec}
\setmainfont{Hiragino Mincho Pro}
\setsansfont{Hiragino Sans}
\setmonofont{Menlo}
\setCJKmainfont{Hiragino Mincho Pro}
\setCJKsansfont{Hiragino Sans}
\setCJKmonofont{Menlo}
% ページ設定
\geometry{margin=1.5cm}
\pagestyle{fancy}
\fancyhf{}
\fancyhead[L]{プログラミング応用 第二回}
\fancyhead[R]{演習問題解説}
\fancyfoot[R]{\ifnum\value{page}=1\small（裏面に続く）\fi}
\fancyfoot[C]{\thepage}
\renewcommand{\headrulewidth}{0.4pt}
\renewcommand{\footrulewidth}{0.4pt}

% 問題番号のカウンター
\newcounter{problem}
\setcounter{problem}{0}

% シンプルな問題環境
\newenvironment{problem}[1][]{
    \refstepcounter{problem}
    \vspace{1em}
    \noindent
    \textbf{問\arabic{problem}} #1
    %\vspace{0.5em}
    %\par
}{
    \vspace{1em}
}

% 解答環境
\newenvironment{solution}{
    \vspace{0.5em}
    \noindent
    \textbf{【解答】}
    \vspace{0.3em}
}{
    \vspace{1em}
}

\begin{document}

% ヘッダー情報
\begin{center}
    {\Large \textbf{プログラミング応用 第二回}}\\[0.5em]
    {\large \textbf{演習問題解説}}\\[1em]
\end{center}

\begin{problem}
講義で紹介した硬貨支払い枚数最小化問題についての以下の問いに答えよ。それぞれの小問は独立している。

\begin{itemize}
\item 講義で紹介した硬貨支払い枚数最小化問題に対する貪欲法は、日本円硬貨(1円、5円、10円、50円、100円、500円)に対して最適であった。講義資料で与えた証明は、「1円玉を5枚以上使えばそれらを5円玉に置き換えて枚数を減らせるため1円玉使う枚数は高々4枚である」という議論を適用して、5円玉は高々1枚、10円玉は高々4枚、50円玉は高々1枚、100円玉は高々4枚であることを用いた。この事実は、\textbf{各硬貨はその一つ上の価値単位の約数になっている}ことによっている（例えば5は10の約数、10は50の約数、50は100の約数）。では、紙幣も追加した場合（1円、5円、10円、100円、500円、1000円、2000円、5000円、10000円）でも貪欲法は最適解を出力するだろうか？正しいならば証明し、そうでないならば反例を示せ。
\item 仮想的な貨幣体系として1円玉、3円玉、4円玉のみからなる体系を考える。この体系では、貪欲法は最適解を出力するだろうか？正しいならば証明し、そうでないならば反例を示せ。
\item 仮想的な通貨体系として1円玉および全ての素数$p$に対して$p$円玉が存在する体系を考える。この体系では、支払う金額 $M \in \mathbb{N}$ がどのような値であっても常に3枚以下の硬貨で支払えることを証明せよ。ただし、ゴールドバッハ予想\footnote{任意の $2$ より大きい偶数は2つの素数の和で表せるという数学の予想。}を仮定してよい。
\end{itemize}
\end{problem}

\begin{solution}

\begin{enumerate}
\item 正しい. 証明は後述.
\item 誤り. $M=6$のとき, 貪欲法の出力は$4$円玉$1$枚と$1$円玉$2$枚であるが, 最適解は$3$円玉$2$枚である.
\item $M=1$のとき, $1$円玉を払えばよい. $M$が偶数のとき, ゴールドバッハ予想より, 二つの素数の和で表せるので, それらに対応する金額を払えばよい. $M\ge 3$が奇数のとき, まず$1$円玉を払った後に偶数のケースが適用できる.
\end{enumerate}

1の証明を以下に記述する. $M$円を支払うとき, 枚数が最小となるときの支払い方の一つを$(k^*_1,k^*_5,k^*_{10},k^*_{50},k^*_{100},k^*_{500}, k^*_{1000}, k^*_{2000}, k^*_{5000}, k^*_{10000})$とする ($k^*_i$は支払う$i$円の枚数を表す).

\begin{itemize}
\item $k^*_1\ge 5$とすると, 1円玉を5枚以上使っており, これらを5円玉に置き換えることでより少ない枚数で支払えてしまい, 最適解であることに矛盾する. 従って $k^*_1\le 4$.
\item 同様の議論で $k^*_5\le 1,k^*_{10}\le 4,k^*_{50}\le 1,k^*_{100}\le 4,k^*_{500}\le 1, k^*_{1000}\le 1$ となる (これは, 一つ上の金額の単位がその一つ下の金額の単位の約数になっていることによる). また, $k^*_{5000} \le 1$ も得られる. すなわち
\begin{align}
  &k^*_1+5k^*_5+10k^*_{10}+50k^*_{50}+100k^*_{100}+500k^*_{500}\le 999, \label{eq:1} \\
  &k^*_1+5k^*_5+10k^*_{10}+50k^*_{50}+100k^*_{100}+500k^*_{500} + 1000k^*_{1000} \le 1999, \label{eq:2} \\
  &5000 k^*_{5000} \le 5000. \label{eq:3}
\end{align}
\item ところが, 2000円の場合, 5000が2000を倍数ではないため, 同様の議論ができない. これをどうするかがポイントとなる.
\begin{enumerate}
  \item ここで, $k^*_{2000}\ge 3$とすると, それだけで6000円を支払っていることになるが, この分を1000円札と5000円札に置き換えることによって, より少ない枚数で支払うことが可能となり, $(k^*_i)$が最適解であることに矛盾する. 従って, $k^*_{2000} \le 2$が成り立つ.
  \item さらに, $(k^*_{1000},k^*_{2000})=(1,2)$とすると, これらを5000円札に置き換えることができる
  \item (a)(b)の議論より以下を得る:
  \begin{align}
    1000k^*_{1000}+2000k^*_{2000} \le 4000. \label{eq:4}
  \end{align}
\end{enumerate}
\item (\ref{eq:1}), (\ref{eq:3}), (\ref{eq:4})より, 1円から5000円を用いて支払う金額は
\begin{align*}
  &k^*_1+5k^*_5+10k^*_{10}+50k^*_{50}+100k^*_{100}+500k^*_{500}+1000k^*_{1000}+2000k^*_{2000}+5000k^*_{5000} \le 9999
\end{align*}
を満たす. よって, 10000円札を使う枚数は$k^*_{10000} = \lfloor M/10000 \rfloor$である. これは貪欲法の出力と一致する.
\item 残った金額を$M' = M - k^*_{10000} \times 10000 \le 9999$とする. 
\item (\ref{eq:1}), (\ref{eq:2})より, 1円から2000円を用いて支払う金額は高々$4999$円なので, 5000円札を使う枚数は$\lfloor M'/5000 \rfloor$である. これも貪欲法の出力と一致する.
\item 残った金額を$M'' = M' - k^*_{5000} \times 5000 \le 4999$とする.
\item (\ref{eq:2})より, 1円から1000円までの金額を用いて支払う金額は高々$1999$円なので, 2000円札を使う枚数は$\lfloor M''/2000 \rfloor$である. これも貪欲法の出力と一致する.
\item 残った金額$M''' = M'' - k^*_{2000} \times 2000 \le 1999$に対して, 講義資料と同様の方法で, 最適解が貪欲法の出力と一致することが示せる.
\end{itemize}

\end{solution}

\begin{problem}
それぞれ$n$個の整数からなる四つの集合$A_1,A_2,A_3,A_4\subseteq \mathbb{Z}$ が与えられる。
ある$a_1\in A_1,a_2\in A_2,a_3\in A_3,a_4\in A_4$ が存在して, $a_1+a_2+a_3+a_4=0$ を満たすかどうかを$O(n^2\log n)$時間で判定するアルゴリズムの疑似コードを記述せよ。
また、そのアルゴリズムの計算量が実際に$O(n^2\log n)$時間であることを簡単に説明せよ（例えばソートの計算量など講義資料で説明した内容は仮定してよい）。
\end{problem}

\begin{solution}
以下のアルゴリズムでよい:
\begin{enumerate}
\item $A_{1,2} = \{a_1+a_2 \mid a_1\in A_1,a_2\in A_2\}$ 計算し, 昇順にソートする ($O(n^2\log n)$時間).
\item 各 $a_3\in A_3, a_4\in A_4$ に対し, $-(a_3+a_4)$ が $A_{1,2}$ に含まれるか二分探索 ($O(n^2\log n)$時間).
\item 含まれていれば, $a_1+a_2+a_3+a_4=0$ を満たす $a_1\in A_1,a_2\in A_2,a_3\in A_3,a_4\in A_4$ が存在するので, 「Yes」を出力して終了.
\item 最後まで終了しなければ, 「No」を出力して終了.
\end{enumerate}
\end{solution}

\newpage

\begin{problem}
凸関数に関する以下のそれぞれの主張は正しいだろうか？正しいならば\textbf{定義に基づいて}証明し、そうでないならば反例を示せ。
\begin{itemize}
\item 関数$f(x)=x^2$は凸関数である。
\item 関数$f(x)=x^3$は凸関数である。
\item 二つの関数$f,g\colon\mathbb{R}^n\to\mathbb{R}$が凸関数であるならば、それらの和$f(x)+g(x)$も凸関数である。
\item 二つの関数$f,g\colon\mathbb{R}^n\to\mathbb{R}$が凸関数であるならば、それらの積$f(x)g(x)$も凸関数である。
\item 二つの関数$f,g\colon\mathbb{R}^n\to\mathbb{R}$が凸関数であるならば、それらの大きい方の値を返す関数$h(x)=\max(f(x),g(x))$も凸関数である。
\end{itemize}

\end{problem}

\begin{solution}
採点は答えだけ確認すればよい (証明はおわなくてよい)
\begin{enumerate}
\item 正しい.
\item 誤り.
\item 正しい.
\item 誤り.
\item 正しい.
\end{enumerate}
\end{solution}

\begin{problem}
関数 $f(x) = x^2-S$ を考える. ただし $S>0$ である. 以下の問いに答えよ.
\begin{enumerate}
\item 点 $(a,f(a))$ における $f$ の接線と $x$軸との交点の座標を求めよ.
\item 数列 $(x_0,x_1,x_2,\dots)$ を以下のように定義する:
\begin{align*}
    x_t = \begin{cases}    
        a & \text{if } t = 0,\\
        \frac{1}{2}\left(x_{t-1} + \frac{S}{x_{t-1}}\right) & \text{if } t \ge 1.
    \end{cases}
\end{align*}
このとき, $\lim_{t\to\infty} x_t$ を求めよ. ただし, $a\ne 0$ とする.
\end{enumerate}


\end{problem}

\begin{solution}
\begin{enumerate}
\item 直線の方程式は $y-a^2 = 2a(x-a)$, すなわち \[ y = 2ax - a^2. \] $a\ne 0$ のとき, $x$軸との交点は $x=\frac{1}{2}\left(a + \frac{S}{a}\right).$ $a=0$のときは, $x$軸上の全ての点が交点となる.
\item $a>0$ならば$\sqrt{S}$, $a<0$ならば$-\sqrt{S}$である.
\end{enumerate}
\end{solution}


\end{document}
